% =============================================================
%  APPENDIX A: AI Acknowledgement
%  ASSIGNED TO: [Person Name]
%  (Include ONLY if AI tools were used for coding or debugging)
% =============================================================

\section{AI Acknowledgement}
\label{app:ai}

% Per the course AI Co-Pilot Rule, if any AI tools were used, this appendix
% must contain:
%   (a) The exact prompts used to query the AI.
%   (b) The raw code generated by the AI.
%   (c) A paragraph explaining the modifications made to ensure physical correctness.
During the creation of this project, AI-based tools (such as ChatGPT, Gemini, and Claude) were used as supporting resources to improve the clarity of explanations, refine LaTeX formatting, and understand certain theoretical concepts that were relevant to this project, along with reading the reference material. We also used these tools to generate our test cases to verify the correctness of our modules' respective implementations.
The contributors facilitated the implementations, analyses, interpretations, and conclusions presented in this project. AI tools served as an integral learning aid and writing resource throughout the process.

\subsection{Prompts Used}

% List each prompt exactly as typed, e.g.:
% \begin{enumerate}
%     \item ``Write a Python function to calculate the Lennard-Jones force between two atoms in reduced units.''
%     \item ``How do I implement the Minimum Image Convention in NumPy?''
% \end{enumerate}

\subsection{AI-Generated Code}

% Include the raw code produced by the AI tool:
% \begin{lstlisting}[caption={Raw AI-generated code for [function name].}]
% [PASTE CODE HERE]
% \end{lstlisting}

\subsection{Modifications and Physical Corrections}

% Write a paragraph describing:
%  - What was wrong or incomplete in the AI output.
%  - What physical constraints were missing (e.g., shift at cutoff, MIC, normalisation).
%  - What changes were made to fix it.

% [WRITE HERE]
